\SourceLeaf{027}
\[
\begin{aligned}
(V-L)\BareRoot{p}&=v,\quad(V'-L')\BareRoot{p'}=v',\\
(V''-L'')\BareRoot{p''}&=v''\text{ etc.}
\end{aligned}
\]
Manifesto itaque res perinde se habebit, ac si observationes immediatae, aequali praecisione gaudentes, puta quarum error medius \(=m\BareRoot{p}=m'\BareRoot{p'}=m''\BareRoot{p''}\) etc., sive quibus pondus \(=1\) tribuitur, suppeditavissent
\[
v=0,\quad v'=0,\quad v''=0\text{ etc.}
\]

\Article{20}
\textsc{Problema.} \textit{Designantibus \(v,v',v''\) etc. functiones lineares indeterminatarum \(x,y,z\) etc. sequentes}
\[
\left.\begin{aligned}
v&=ax+by+cz+\text{etc.}+l\\
v'&=a'x+b'y+c'z+\text{etc.}+l'\\
v''&=a''x+b''y+c''z+\text{etc.}+l''\text{ etc.}
\end{aligned}\right\}\quad\text{(I)}
\]
\textit{ex omnibus systematibus coëfficientium \(\kappa,\kappa',\kappa''\) etc., qui indefinite dant}
\[
\kappa v+\kappa'v'+\kappa''v''+\text{etc.}=x-k
\]
\textit{ita ut \(k\) sit quantitas determinata i. e. ab \(x,y,z\) etc. independens, eruere id, pro quo \(\kappa\kappa+\kappa'\kappa'+\kappa''\kappa''+\text{etc.}\) nanciscatur valorem minimum.}

\textit{Solutio.} Statuamus
\[
\left.\begin{aligned}
av+a'v'+a''v''+\text{etc.}&=\xi\\
bv+b'v'+b''v''+\text{etc.}&=\eta\\
cv+c'v'+c''v''+\text{etc.}&=\zeta
\end{aligned}\right\}\quad\text{(II)}
\]
etc.: eruntque etiam \(\xi,\eta,\zeta\) etc. functiones lineares ipsarum \(x,y,z\) etc., puta
\[
\left.\begin{aligned}
\xi&=x\Sigma aa+y\Sigma ab+z\Sigma ac+\text{etc.}+\Sigma al\\
\eta&=x\Sigma ab+y\Sigma bb+z\Sigma bc+\text{etc.}+\Sigma bl\\
\zeta&=x\Sigma ac+y\Sigma bc+z\Sigma cc+\text{etc.}+\Sigma cl\text{ etc.}
\end{aligned}\right\}\quad\text{(III)}
\]
(ubi \(\Sigma aa\) denotat aggregatum \(aa+a'a'+a''a''+\text{etc.}\), ac perinde de reliquis) multitudoque ipsarum \(\xi,\eta,\zeta\) etc. multitudini indeterminatarum \(x,y,z\) etc. aequalis, puta \(=\rho\). Per eliminationem itaque elici poterit aequatio talis\footnote{Ratio, cur ad denotandos coëfficientes e tali eliminatione prodeuntes, hos potissimum characteres elegerimus, infra elucebit.}
\SourceLeaf{028}
\[
x=A+[\alpha\alpha]\xi+[\alpha\beta]\eta+[\alpha\gamma]\zeta+\text{etc.}
\]
in qua substituendo pro \(\xi,\eta,\zeta\) etc. valores earum ex III, aequatio identica prodire debet. Quare statuendo
\[
\left.\begin{aligned}
a[\alpha\alpha]+b[\alpha\beta]+c[\alpha\gamma]+\text{etc.}&=\alpha\\
a'[\alpha\alpha]+b'[\alpha\beta]+c'[\alpha\gamma]+\text{etc.}&=\alpha'\\
a''[\alpha\alpha]+b''[\alpha\beta]+c''[\alpha\gamma]+\text{etc.}&=\alpha''\text{ etc.}
\end{aligned}\right\}\quad\text{(IV)}
\]
necessario erit, indefinite
\[
\alpha v+\alpha'v'+\alpha''v''+\text{etc.}=x-A\quad\text{(V)}
\]
Haec aequatio docet, inter systemata valorum coëfficientium \(\kappa,\kappa',\kappa''\) etc. certo etiam referendos esse hos \(\kappa=\alpha,\kappa'=\alpha',\kappa''=\alpha''\) etc., nec non, pro systemate quocunque, fieri debere indefinite
\[
(\kappa-\alpha)v+(\kappa'-\alpha')v'+(\kappa''-\alpha'')v''+\text{etc.}=A-k
\]
quae aequatio implicat sequentes
\[
\begin{aligned}
(\kappa-\alpha)a+(\kappa'-\alpha')a'+(\kappa''-\alpha'')a''+\text{etc.}&=0\\
(\kappa-\alpha)b+(\kappa'-\alpha')b'+(\kappa''-\alpha'')b''+\text{etc.}&=0\\
(\kappa-\alpha)c+(\kappa'-\alpha')c'+(\kappa''-\alpha'')c''+\text{etc.}&=0\text{ etc.}
\end{aligned}
\]
Multiplicando has aequationes resp. per \([\alpha\alpha],[\alpha\beta],[\alpha\gamma]\) etc., et addendo, obtinemus propter (IV)
\[
(\kappa-\alpha)\alpha+(\kappa'-\alpha')\alpha'+(\kappa''-\alpha'')\alpha''+\text{etc.}=0
\]
sive quod idem est
\[
\begin{aligned}
&\kappa\kappa+\kappa'\kappa'+\kappa''\kappa''+\text{etc.}\\
&=\alpha\alpha+\alpha'\alpha'+\alpha''\alpha''+\text{etc.}{}\\
&\quad+(\kappa-\alpha)^2+(\kappa'-\alpha')^2+(\kappa''-\alpha'')^2+\text{etc.}
\end{aligned}
\]
unde patet, aggregatum \(\kappa\kappa+\kappa'\kappa'+\kappa''\kappa''+\text{etc.}\) valorem minimum obtinere, si statuatur \(\kappa=\alpha,\kappa'=\alpha',\kappa''=\alpha''\) etc. Q.~E.~I.

Ceterum hic valor minimus ipse sequenti modo eruitur. Aequatio (V) docet, esse
\SourceLeaf{029}
\[
\begin{aligned}
\alpha a+\alpha'a'+\alpha''a''+\text{etc.}&=1\\
\alpha b+\alpha'b'+\alpha''b''+\text{etc.}&=0\\
\alpha c+\alpha'c'+\alpha''c''+\text{etc.}&=0\text{ etc.}
\end{aligned}
\]
Multiplicando has aequationes resp. per \([\alpha\alpha],[\alpha\beta],[\alpha\gamma]\) etc. et addendo, protinus habemus adiumento aequationum (IV)
\[
\alpha\alpha+\alpha'\alpha'+\alpha''\alpha''+\text{etc.}=[\alpha\alpha]
\]

\Article{21}
Quum observationes suppeditaverint aequationes (proxime veras) \(v=0,v'=0,v''=0\) etc., ad valorem incognitae \(x\) inde eliciendum combinatio illarum aequationum talis
\[
\kappa v+\kappa'v'+\kappa''v''+\text{etc.}=0
\]
adhibenda est, quae ipsi \(x\) coëfficientem \(1\) conciliet, incognitasque reliquas \(y,z\) etc. eliminet; cui determinationi per art. 18 pondus
\[
=\frac1{\kappa\kappa+\kappa'\kappa'+\kappa''\kappa''+\text{etc.}}
\]
tribuendum erit. Ex art. praec. itaque sequitur, determinationem maxime idoneam eam fore, ubi statuatur \(\kappa=\alpha,\kappa'=\alpha',\kappa''=\alpha''\) etc. Hoc pacto \(x\) obtinet valorem \(A\), manifestoque idem valor etiam (absque cognitione multiplicatorum \(\alpha,\alpha',\alpha''\) etc.) protinus per eliminationem ex aequationibus \(\xi=0,\eta=0,\zeta=0\) etc. elici potest. Pondus huic determinationi tribuendum erit \(=\frac1{[\alpha\alpha]}\), sive error medius in ipsa metuendus
\[
=m\BareRoot{p[\alpha\alpha]}=m'\BareRoot{p'[\alpha\alpha]}=m''\BareRoot{p''[\alpha\alpha]}\text{ etc.}
\]

Prorsus simili modo determinatio maxime idonea incognitarum reliquarum \(y,z\) etc. eosdem valores ipsis conciliabit, qui per eliminationem ex iisdem aequationibus \(\xi=0,\eta=0,\zeta=0\) etc. prodeunt.

Denotando aggregatum indefinitum \(vv+v'v'+v''v''\text{ etc.}\), sive quod idem est hoc
\[
p(V-L)^2+p'(V'-L')^2+p''(V''-L'')^2+\text{etc.}
\]
